\documentclass[12pt]{article}
\usepackage[utf8]{inputenc}
\usepackage[T1]{fontenc}
\usepackage{hyperref}
\title{Traveller Scottish Language Resources for K-12 Curriculum Development}
\author{Compiled by Claude Sonnet 5 (claude-sonnet-5) for Andr\'as Kornai}
\date{}
\begin{document}
\maketitle

\section{Introduction}
Traveller Scottish (Scottish Cant; ISO 639-3 \texttt{trl}) is a para-Romani or mixed language of Scottish Lowland Travellers. Reports give about 460 native speakers (2011) and describe a largely oral tradition with local dialects. Estimates say up to half of the vocabulary is Romani-derived, with many Scots words. Highland Travellers are a separate group and should not be confused with it. Only short phrase and word lists were found.

\section{Primary Educational Resources}
\begin{itemize}
\item Language cards with greetings, counting and phrases (Gypsy Roma Traveller History Month 2022), with translations credited in the cards to Davie Donaldson: \url{https://www.gypsy-traveller.org/wp-content/uploads/2022/05/A_FFT_GRTHM_2022_Language-Cards_a4-1.pdf}
\item A copy of the language cards hosted by Telford council: \url{https://www.telford.gov.uk/download/downloads/id/21126/gypsy_roma_traveller_history_month_-_language_cards.pdf}
\end{itemize}

\section{Linguistic and Reference Materials}
\begin{itemize}
\item ``Multi-Layered Communication and Function in Scottish Traveller Cant'', \emph{Traditiones} (2021): \url{https://ojs.zrc-sazu.si/traditiones/article/view/10534}
\item Full text of the same article: \url{https://ojs.zrc-sazu.si/traditiones/article/download/10534/9696/29588}
\item Cant glossary on travellers.scot: \url{https://travellers.scot/?p=91}
\item Overview page: \url{https://www.gypsy-traveller.org/scottish-cant-or-scots-romani/}
\end{itemize}

\section{Religious and Community Resources}
The Gypsy Traveller organisation page above is the only community source found. No scripture was found.

\section{Digital and Audio Resources}
None were found in the search results.

\section{Curriculum Development Context}
The resources found are short phrase cards and a glossary, which suit awareness or heritage lessons and do not support a full curriculum. Cant is traditionally a covert oral language, so Traveller communities would have to decide what is shared and how. Kirk and \'O Baoill (2002), \emph{Travellers and their Language}, and Clement (1981) are cited in the Traditiones article. I did not find them online.
\end{document}